\chapter{\texorpdfstring{\(\mathcal N=4\)}{N=4} Super-Yang--Mills as a Conformal Theory}
\label{ch:volume-iii-n4-super-yang-mills-conformal-theory}

\section{Field Content and Symmetries}

Let \(G\) be a compact simple Lie group with Lie algebra \(\mathfrak g\).
Let \(\operatorname{Tr}\) be an invariant positive quadratic form on
\(\mathfrak g\), normalized for \(G=SU(N)\) by
\[
  \operatorname{Tr}(t^at^b)=\frac12\delta^{ab}
\]
in the fundamental representation.  Four-dimensional
\(\mathcal N=4\) super-Yang--Mills has the following adjoint-valued fields:
\[
  A_\mu,\qquad
  \lambda^I_\alpha,\qquad
  \bar\lambda_{\dot\alpha I},\qquad
  \phi^m.
\]
Here \(A_\mu\) is a gauge connection, \(\lambda^I_\alpha\) and
\(\bar\lambda_{\dot\alpha I}\) are Weyl fermions, \(I=1,\ldots,4\), and
\(\phi^m\), \(m=1,\ldots,6\), are real scalar fields.  All fields transform in
the adjoint representation of \(G\).  The index \(m\) transforms in the vector
representation \(\mathbf6\) of \(SO(6)_R\simeq SU(4)_R\), while \(I\) is an
\(SU(4)_R\) fundamental index.

The curvature and covariant derivative are
\[
  F_{\mu\nu}
  =
  \partial_\mu A_\nu-\partial_\nu A_\mu-\ii[A_\mu,A_\nu],
  \qquad
  D_\mu X=\partial_\mu X-\ii[A_\mu,X]
\]
for any adjoint-valued field \(X\).  With Lorentzian signature, the action has
the form
\begin{equation}
\begin{aligned}
  S_{\mathcal N=4}
  =
  \frac1{g_{\rm YM}^2}\int\dd^4x\,\operatorname{Tr}\bigg[
  &-\frac12F_{\mu\nu}F^{\mu\nu}
  -D_\mu\phi^mD^\mu\phi^m
  +\ii\bar\lambda_I\bar\sigma^\mu D_\mu\lambda^I\\
  &+\frac12[\phi^m,\phi^n][\phi^m,\phi^n]
  +\rho^m_{IJ}\lambda^I[\phi^m,\lambda^J]
  +\bar\rho^{m\,IJ}\bar\lambda_I[\phi^m,\bar\lambda_J]
  \bigg].
\end{aligned}
  \label{eq:n4-sym-action}
\end{equation}
The tensors \(\rho^m_{IJ}=-\rho^m_{JI}\) are Clebsch--Gordan coefficients for
the map
\[
  \mathbf4\wedge\mathbf4\longrightarrow\mathbf6
\]
of \(SU(4)_R\), and \(\bar\rho^{m\,IJ}\) are their conjugates.  Supersymmetry
fixes the relative coefficients in \eqref{eq:n4-sym-action}.  Thus the gauge
coupling, Yukawa coupling, and scalar quartic coupling are a single coupling.

\section{One-Loop Beta Function}

For a general four-dimensional gauge theory with \(n_f\) Weyl fermions and
\(n_s\) real scalars in the adjoint representation, the one-loop beta function
for \(G=SU(N)\) is
\begin{equation}
  \mu\frac{\dd g_{\rm YM}}{\dd\mu}
  =
  -\frac{g_{\rm YM}^3}{(4\pi)^2}
  \frac N3
  \left(11-2n_f-\frac12 n_s\right)
  +O(g_{\rm YM}^5).
  \label{eq:adjoint-gauge-theory-one-loop-beta}
\end{equation}
The \(\mathcal N=4\) multiplet has
\[
  n_f=4,
  \qquad
  n_s=6.
\]
Substitution into \eqref{eq:adjoint-gauge-theory-one-loop-beta} gives
\[
  11-2\cdot4-\frac12\cdot6=0.
\]
The perturbative trace anomaly associated with the running gauge coupling
therefore has zero one-loop coefficient.

Maximal supersymmetry strengthens this cancellation.  The exactly marginal
parameter is the complex coupling
\begin{equation}
  \tau_{\rm YM}
  =
  \frac{\theta}{2\pi}
  +
  \frac{4\pi\ii}{g_{\rm YM}^2},
  \label{eq:n4-complex-coupling}
\end{equation}
where \(\theta\) is the coefficient of
\[
  \frac1{8\pi^2}\int \operatorname{Tr}(F\wedge F).
\]
The supersymmetric non-renormalization theorem for the \(\mathcal N=4\)
stress-tensor multiplet implies
\[
  \beta_{\tau_{\rm YM}}=0.
\]
Thus each value of \(\tau_{\rm YM}\) defines a conformal theory.  The
superconformal algebra throughout this conformal manifold is
\[
  \mathfrak{psu}(2,2|4).
\]

\section{Protected Scalar Operators}

The six scalars \(\phi^m\) transform as the \(\mathbf6=[0,1,0]\) of
\(SU(4)_R\).  The traceless symmetric bilinear
\begin{equation}
  \mathcal O^{mn}
  =
  \operatorname{Tr}(\phi^m\phi^n)
  -
  \frac16\delta^{mn}
  \operatorname{Tr}(\phi^k\phi^k)
  \label{eq:n4-stress-tensor-multiplet-primary}
\end{equation}
transforms in the representation
\[
  \mathbf{20}'=[0,2,0].
\]
It is a scalar superconformal primary.  Since it belongs to the half-BPS
family \([0,p,0]\) with \(p=2\), its dimension is fixed by
\eqref{eq:half-bps-bound}:
\[
  \Delta(\mathcal O^{mn})=2.
\]
This operator is the primary of the stress-tensor multiplet.  The conserved
stress tensor, the \(SU(4)_R\) currents, and the supercurrents are descendants
inside the same shortened multiplet.

The singlet
\[
  \mathcal K=\operatorname{Tr}(\phi^k\phi^k)
\]
has the same classical dimension but belongs to a long multiplet.  Its
dimension is governed by dynamical anomalous dimensions.  It is the Konishi
primary, and its scaling dimension is a nontrivial function of
\(\tau_{\rm YM}\).

\begin{figure}[H]
  \centering
  \resizebox{0.92\textwidth}{!}{%
  \begin{tikzpicture}[x=1cm,y=1cm,every node/.style={font=\scriptsize}]
    \tikzset{
      arrow/.style={-{Latex[length=2mm]},line width=0.85pt},
      box/.style={rounded corners=4pt,draw,line width=0.8pt,fill=blue!4,
        minimum width=3.1cm,minimum height=1.0cm}
    }
    \node[box] (fields) at (-4.2,0)
      {\(\phi^m\in\mathbf6=[0,1,0]\)};
    \node[box] (bilinear) at (0,0.8)
      {\(\mathcal O^{mn}\in[0,2,0]\)};
    \node[box] (konishi) at (0,-0.8)
      {\(\mathcal K\in[0,0,0]\)};
    \node[box] (protected) at (4.3,0.8)
      {half-BPS: \(\Delta=2\)};
    \node[box] (long) at (4.3,-0.8)
      {long: \(\Delta\) varies};
    \draw[arrow] (fields)--(bilinear);
    \node[anchor=east] at (-1.72,1.36) {traceless symmetric};
    \draw[arrow] (fields)--(konishi)
      node[midway,below] {singlet};
    \draw[arrow,green!45!black] (bilinear)--(protected);
    \draw[arrow,orange!85!black] (konishi)--(long);
    \node[align=center,violet!80!black] at (0,1.85)
      {same classical dimension, different superconformal representation};
  \end{tikzpicture}
  }
  \caption{The protected stress-tensor multiplet primary and the Konishi
  primary are distinguished by their \(SU(4)_R\) representation and by
  shortening.}
  \label{fig:n4-protected-versus-konishi}
\end{figure}

\section{Local Operators as Conformal Data}

At a conformal point, local gauge-invariant operators are organized by
irreducible representations of \(\mathfrak{psu}(2,2|4)\).  A convenient
generating class is provided by traces of adjoint fields and covariant
derivatives:
\[
  \operatorname{Tr}
  \left(
    X_1X_2\cdots X_L
  \right),
  \qquad
  X_\ell\in
  \{F_{\mu\nu},\,D_{\mu_1}\cdots D_{\mu_k}\phi^m,\,
    D_{\mu_1}\cdots D_{\mu_k}\lambda^I_\alpha,\ldots\}.
\]
The integer \(L\) is the trace length before operator mixing is diagonalized.
At finite \(N\), trace identities and multi-trace mixing are part of the
operator algebra.  At large \(N\), the single-trace sector has a simpler
leading structure, developed in the next chapters.
